% sec-10-relevant-information.tex - Relevant Information (full stage)
\section{Relevant Information}

This section provides the additional relevant information that the U.S. Food and
Drug Administration may request under 21 CFR \S 312.23(a)(11), together with the
method-level grounding and verification disclosures that apply because this
Investigational New Drug application (IND v1.0) was assembled by an artificial
intelligence build operating over the repository at version v4.3.0. The
disclosures are placed here, at the close of the package, so that a reviewer can
read the entire submission first and then evaluate exactly how the document was
constructed, how each figure and number was grounded, how the build is monitored
in real time, and where the residual limitations lie. Nothing in this section
substitutes for the clinical, statistical, and regulatory substance of the
preceding sections; rather, it states the controls that give a reviewer
confidence that the substance is internally consistent, traceable to named source
files, and reproducible. The verification posture described here is aligned with
the verification-before-generation discipline that the program has advanced for
physical artificial intelligence in oncology trials \cite{congress9510,
clintrustpai}, and it is consistent with the phased-development expectations for
an initial Phase 1 submission \cite{phase1guidance}.

\subsection{Method-Level Grounding and Verification}

The IND was produced by a single-prompt, multi-file build that proceeds through
four stages, mermaid to draft to full to final, with the present full stage
expanding each draft scaffold into regulator-grade prose, tables, and figures.
To keep that build verifiable rather than merely fluent, the method applies five
name-matching and grounding verifications to every artifact it emits. The five
verifications are summarized in Figure 23 and detailed, with a concrete IND
example for each, in Table 17. They are not abstract quality slogans; each one is
a check that a reviewer can independently reproduce by inspecting the public
repository tree, the rendered document, and the commit and pull-request history.
The five together yield a grounded, monitorable IND in which every figure names
its source, every figure is cited where its section discusses it, every file is
visible on GitHub as it is written, the repository and directory names mirror the
document structure, and the section file names map one to one onto the regulatory
sections they contain.

\begin{mermaidfig}
\node[mmdec] (k1) at (0,0) {1 Grounding};
\node[mmdec] (k2) at (0,-2.0) {2 Figures match context};
\node[mmdec] (k3) at (0,-4.0) {3 GitHub real-time};
\node[mmdec] (k4) at (0,-6.0) {4 Repo / dir names match};
\node[mmdec] (k5) at (0,-8.0) {5 File names match};
\node[mmin] (e1) at (5.0,0) {22 grayscale Mermaid figures, each sourced};
\node[mmin] (e2) at (5.0,-2.0) {each figure cited where its \S section discusses it};
\node[mmin] (e3) at (5.0,-4.0) {one commit per file, auto-pushed to one PR};
\node[mmin] (e4) at (5.0,-6.0) {trial-ind/ matches the repository tree};
\node[mmin] (e5) at (5.0,-8.0) {sec-06-cmc.tex matches the CMC section};
\node[mmgoal] (g) at (10.4,-4.0) {Grounded,\\monitorable IND};
\draw[mmedge] (k1) -- (e1);
\draw[mmedge] (k2) -- (e2);
\draw[mmedge] (k3) -- (e3);
\draw[mmedge] (k4) -- (e4);
\draw[mmedge] (k5) -- (e5);
\draw[mmedge] (e1) to[out=0,in=150,looseness=0.6] (g.west);
\draw[mmedge] (e2) to[out=0,in=165,looseness=0.6] (g.west);
\draw[mmedge] (e3) -- (g.west);
\draw[mmedge] (e4) to[out=0,in=195,looseness=0.6] (g.west);
\draw[mmedge] (e5) to[out=0,in=210,looseness=0.6] (g.west);
\end{mermaidfig}
\figcaption{Figure 23. Five name-matching and grounding verifications, each with a
concrete IND example: grounding (every one of the 22 grayscale figures names its
source files); figures match context (each is cited where its section discusses
it); GitHub real-time (one commit per file, auto-pushed to one continuously
updated pull request); repository and directory names match (the trial-ind tree
mirrors the document structure); and file names match (sec-06-cmc.tex maps to the
Chemistry, Manufacturing and Control section). Together they yield a grounded,
monitorable IND.}

\subsubsection{The Five Verifications in Detail}

\textit{Verification 1, grounding.} Every quantitative claim that a figure carries
is traceable to a named source file rather than invented at generation time. The
22 grayscale figures of this IND each declare, in their Mermaid source under
\texttt{trial-ind/mermaid/}, the source files from which their nodes, edge labels,
and numbers were drawn, and those numbers are the same document counts, dose
levels, force caps, sampling rates, and review clocks reused in the surrounding
prose. For example, the dose levels of 160 mg, 220 mg, and 300 mg once daily, the
28 day dose-limiting-toxicity window, and the maximum enrollment of 18 treated
subjects appear identically in the dose-escalation figure, in the General
Investigational Plan, and in the protocol from which they originate. This binds
the figures to the trial facts and prevents a figure from drifting away from the
text it illustrates \cite{oncotrialllm, pdacdigtwinp}.

\textit{Verification 2, figures match context.} Each figure is cited in the
sentence or paragraph that discusses it, in the section where that discussion
belongs, so that no figure floats unreferenced and no claim in the prose lacks its
illustrating figure. The mechanism-of-action figure is cited in the introductory
material that names the drug and its pharmacological class, the 3+3 automaton is
cited in the General Investigational Plan, the chemistry information architecture
is cited in the Chemistry, Manufacturing and Control section, and the two figures
of the present section are cited here. A reviewer can confirm this by following any
figure number from its caption back to the citing sentence.

\textit{Verification 3, GitHub real time.} Each distinguishable file is committed
and pushed at the moment it is written, and those commits accumulate on a single
continuously updated pull request, so that the IND is observable on GitHub as it is
constructed rather than only after completion. This gives the sponsor-investigator,
and in principle an auditor, a real-time view of provenance: who or what wrote each
file, in what order, and against which commit. The auto-commit and auto-pull-request
behavior is the documentary analogue of the audit-trail and monitorability
expectations that govern the device-directed component of the trial \cite{paiautospons,
paisitedocpk}.

\textit{Verification 4, repository and directory names match.} The directory tree
named in the prose is the directory tree on disk. References to
\texttt{trial-ind/mermaid/}, \texttt{trial-ind/full-ind/sections/}, and
\texttt{trial-ind/inputs/} denote the actual directories of the repository at
version v4.3.0, so a reader can navigate from any path mentioned in the document
straight to the corresponding location in the public tree without translation.

\textit{Verification 5, file names match.} The section file names map one to one
onto the regulatory sections they contain. \texttt{sec-06-cmc.tex} is the
Chemistry, Manufacturing and Control section, \texttt{sec-07-pharmacology-toxicology.tex}
is the Pharmacology and Toxicology section, \texttt{sec-10-relevant-information.tex}
is the present section, and \texttt{fig-09-five-grounding-verifications.md} is the
Mermaid source for Figure 23. The naming convention makes the package
self-describing and removes ambiguity about which file holds which content.

\vspace{0.2em}
\noindent\begin{tabularx}{\textwidth}{L{0.6cm} L{4.4cm} Y}
\toprule
\textbf{\#} & \textbf{Verification check} & \textbf{Concrete IND example} \\
\midrule
1 & Grounding via Mermaid figures & The 22 grayscale figures in \texttt{trial-ind/mermaid/} each name their source files and carry the dose levels (160, 220, 300 mg), the 28 day DLT window, and the 30 day FDA clock reused verbatim in the prose \\
2 & Figures match IND context & Each figure is cited in the sentence that discusses it, in the section where that discussion belongs (mechanism in the Introduction, 3+3 in the Plan, CMC architecture in the CMC section, these two figures here) \\
3 & AI files viewable on GitHub in real time & One commit per file is pushed at generation, accumulating on a single continuously updated pull request, so each artifact appears on the branch as it is made \\
4 & Repository and directory names match & The prose refers to \texttt{trial-ind/mermaid} and \texttt{full-ind/sections}, which are the directories on disk in repository v4.3.0 \\
5 & File names match repository file names & A cited \texttt{sec-06-cmc.tex} is the CMC section; \texttt{fig-09-five-grounding-verifications.md} is the Mermaid source for Figure 23 \\
\bottomrule
\end{tabularx}
\figcaption{Table 17. The five name-matching and grounding verifications, each with
a concrete example drawn from this IND. Every check is independently reproducible
from the public repository tree, the rendered document, and the commit and
pull-request history.}

\subsection{Figure Grounding: Mermaid to TikZ}

A second, finer-grained discipline governs the figures themselves: every figure in
this IND is authored once as a grayscale Mermaid source in
\texttt{trial-ind/mermaid/} and then reproduced as a TikZ figure of identical
complexity in the LaTeX document, with no simplification of nodes, edges, edge
labels, or quantitative content. The translation pathway and its verify-twice
checkpoint are shown in Figure 24. The Mermaid source defines the nodes, edges,
grayscale palette, and quantitative data; the conversion preserves all of these;
the TikZ \texttt{mermaidfig} environment renders them using the \texttt{mm*} node
styles defined in \texttt{indstyle.sty} (the same eight-tone ramp that maps one to
one from the Mermaid classDef classes); and the result is verified twice, for
text-box and arrow overlaps, for the specified arrow looseness, and for proper box
spacing, before it appears identically in the draft, full, and final stages. A
figure that fails the verify-twice checkpoint is returned to the TikZ step and
corrected; only a figure that passes is allowed to render in the document.

\begin{mermaidfig}
\node[mmin] (src) at (0,0) {mermaid/fig-NN.md\\(GitHub Mermaid)};
\node[mmproc] (con) at (4.4,0) {nodes, edges, grayscale palette, quantitative data};
\node[mmproc] (tikz) at (8.8,0) {TikZ mermaidfig\\(mm* node styles)};
\node[mmgoal] (out) at (13.2,0) {Identical figure in draft / full / final};
\node[mmdec] (v) at (13.2,-3.0) {Verify twice: no overlaps, arrow looseness, box spacing};
\draw[mmedge] (src) -- (con);
\draw[mmedge] (con) -- (tikz);
\draw[mmedge] (tikz) -- (out);
\draw[mmedge] (out) -- (v);
\draw[mmedge] (v.west) to[out=180,in=-90,looseness=0.8] node[mmlabel,pos=0.55] {pass} (out.south west);
\draw[mmedged] (v.north west) to[out=150,in=-60,looseness=0.7] node[mmlabel,pos=0.5] {fail} (tikz.south);
\end{mermaidfig}
\figcaption{Figure 24. Figure grounding. Each Mermaid source in \texttt{mermaid/} is
reproduced as a TikZ \texttt{mermaidfig} of the same complexity (the same nodes,
edges, grayscale palette, and quantitative data), then verified twice for text-box
and arrow overlaps, for the specified arrow looseness, and for proper box spacing,
before it appears identically in the draft, full, and final stages. A figure that
fails is returned to the TikZ step; only a figure that passes renders.}

The practical effect of this discipline is that the figure a reviewer sees in the
compiled PDF is provably the same diagram, with the same numbers, as the Mermaid
source committed to the repository. There is no separate figure-rendering pipeline,
no raster export, and no opportunity for a number to be transcribed differently
between the source and the rendered figure. All figures are vector TikZ;
the IND contains no PNG or JPG images and no \texttt{\textbackslash includegraphics}
calls. The grayscale ramp is fixed in the style file, so the tonal encoding of node
roles (goal, process, accent, step, input, context, decision, dark) is identical
across all 22 figures and across the four build stages \cite{fdadigtwinpc}.

\subsection{Real-Time Auto-Commit and Auto-Pull-Request Monitorability}

The build is monitorable in real time. As each section file, style file, figure
source, and supporting file is written, it is committed individually and pushed to
the remote, and the accumulating commits are surfaced on a single open pull request
that updates continuously over the life of the build. This means that the
provenance of the IND is not reconstructed after the fact from a final snapshot;
it is visible while the document is being produced. A reader with access to the
public repository can watch the \texttt{trial-ind/full-ind/} tree fill in section
by section, can inspect the diff for any individual file, and can read the commit
history as an ordered, timestamped record of how the package came together. This
real-time monitorability is the documentary counterpart of the trial's own
audit-trail commitments, under which the device preserves a window of records
around each physical artificial intelligence adverse event, and it is consistent
with the broader transparency posture that the program has argued is necessary
before higher-autonomy systems are used in oncology trials \cite{congress9510,
clintrustpai, natlpaiplatf}. The same single-prompt, auto-commit, auto-pull-request
method has been demonstrated across prior multi-file builds in the program,
including the companion Phase 2 protocol and the legislative-text repositories,
which establishes that the approach is reproducible rather than specific to this
submission \cite{oncotrialllm, paipredict4x}.

\subsection{Repository, Directory, and File Name Matching}

The repository, directory, and file naming is itself a control. The IND lives under
\texttt{trial-ind/} in the public repository, and within it the four build stages
occupy \texttt{mermaid/}, \texttt{draft-ind/}, \texttt{full-ind/}, and
\texttt{final-ind/}, mirroring the mermaid to draft to full to final progression.
Within each LaTeX stage the section files follow a fixed
\texttt{sec-NN-<name>.tex} convention whose number and name match the ReGARDD table
of contents, so that the file that holds the Chemistry, Manufacturing and Control
section is named \texttt{sec-06-cmc.tex}, the file that holds the Pharmacology and
Toxicology section is named \texttt{sec-07-pharmacology-toxicology.tex}, and the
figure that this section reproduces first is sourced from
\texttt{fig-09-five-grounding-verifications.md}. Table 18 records the directory and
file-naming correspondences that make the package self-describing, so that a
reviewer can move from any path named in the prose directly to the corresponding
location on disk without translation, and can confirm that no section is missing
and no file is misnamed.

\vspace{0.2em}
\noindent\begin{tabularx}{\textwidth}{L{4.6cm} L{4.2cm} Y}
\toprule
\textbf{Path in repository v4.3.0} & \textbf{Holds} & \textbf{Name-match guarantee} \\
\midrule
\texttt{trial-ind/} & The full IND across four stages & Top-level directory named for the IND it contains \\
\texttt{trial-ind/mermaid/} & The 22 grayscale figure sources & Stage 1 directory; \texttt{fig-NN-<name>.md} files \\
\texttt{trial-ind/full-ind/sections/} & The full-stage section files & Stage 3 sections; one \texttt{sec-NN-<name>.tex} per top-level section \\
\texttt{sec-06-cmc.tex} & Chemistry, Manufacturing and Control & File name matches the CMC section it contains \\
\texttt{sec-10-relevant-information.tex} & Relevant Information (this section) & File name matches the section it contains \\
\texttt{fig-09-five-grounding-verifications.md} & Source for Figure 23 & Figure file name matches the figure it renders \\
\bottomrule
\end{tabularx}
\figcaption{Table 18. Repository, directory, and file name-matching for the IND
(Rule 5). Every path named in the prose denotes the actual location on disk in
repository v4.3.0, making the package self-describing and navigable.}

\subsection{Limitations of the AI-Generation Method}

The grounding and verification controls above constrain the build but do not
eliminate its limitations, and those limitations are stated plainly here so that a
reviewer weighs the package accordingly. Table 19 enumerates the principal
limitations together with the mitigation applied to each. The most important point
is the one that no method-level control can change: the AI-generation method alters
none of the clinical realities of the trial. It does not change tumor biology, it
does not shorten the 28 day dose-limiting-toxicity observation window or the 30 and
90 day safety follow-up or the survival follow-up to 24 months, and it does not
compress the fixed 30 day FDA review clock or the IRB meeting calendar. What the
method compresses is the administrative and document-preparation bucket alone; the
clinical and operational bucket and the fixed regulatory-review bucket are
unchanged \cite{phase1guidance, pdac060s2030}. Every generated document, this IND
included, must still pass human medical, statistical, and regulatory review before
use, and the build is an independent research artifact, not an active IND or
investigational device exemption and not regulatory advice
\cite{cfr312paiuni, cfr050paiuni}.

Several narrower limitations follow from how the build operates. The method does not
rely on live internet search for a project of this size, because specifying single
external URLs is unreliable at scale; it depends instead on a curated repository of
appropriately sized inputs, which bounds the source material a single build can
absorb. Because the model does not always have access to a PDF compiler during
generation, a portion of senior-author formatting, including table column widths,
figure-caption spacing, and occasional page breaks, must still be applied by hand,
and the final stage performs that pass explicitly. Total input is kept under roughly
5 MB with individual files held under per-file token limits, which caps how much
source material a single build can ingest. High-dimensional or approximate inputs
can make outputs less deterministic, and oversized or overly complex inputs can
surface server errors that signal reduced quality; the build mitigates this by
keeping inputs curated and appropriately sized. These constraints are the same ones
characterized for the broader single-prompt method and are disclosed here for the
specific case of an IND \cite{ichpaistduni, chatgpt100kp, llmcomp16fda}.

\vspace{0.2em}
\noindent\begin{tabularx}{\textwidth}{L{4.8cm} Y}
\toprule
\textbf{Limitation} & \textbf{Mitigation applied in this IND} \\
\midrule
The method changes no clinical reality (tumor biology, DLT window, survival follow-up) & The 28 day DLT window, 30 and 90 day safety follow-up, and 24 month survival follow-up are preserved exactly as specified in the protocol; the build compresses only document preparation \\
The method cannot shorten fixed regulatory clocks & The 30 day FDA review and the IRB meeting calendar are stated as binding and unchanged; faster authoring only starts the clock sooner \\
Generated content can be fluent but wrong without grounding & The five name-matching and grounding verifications and the figure verify-twice checkpoint bind every figure and number to a named source file \\
No live internet search at scale; URLs unreliable in bulk & Inputs are a curated repository of appropriately sized files rather than ad hoc external URLs \\
No PDF compiler during generation; some formatting needs a manual pass & Senior-author formatting (column widths, caption spacing, page breaks) is applied explicitly in the final stage \\
Bounded input budget (under ~5 MB total; per-file token limits) & Source material is curated and apportioned so each build stays within limits without truncation \\
High-dimensional or oversized inputs reduce determinism and can surface server errors & Inputs are kept appropriately sized; complex content is split so the build remains deterministic and stable \\
The artifact is research, not an active regulatory submission & The package is labeled an independent research artifact, not an active IND or IDE and not regulatory advice; human review is required before any use \\
\bottomrule
\end{tabularx}
\figcaption{Table 19. Principal limitations of the AI-generation method and the
mitigation applied to each in this IND. The clinical and fixed-regulatory
limitations are intrinsic and cannot be removed; the method-level limitations are
mitigated by grounding, verification, and a human review requirement.}

\subsection{Compilation Statement}

All LaTeX sources for this IND compile in Overleaf with pdfLaTeX. The full-stage
document is built from \texttt{main.tex}, the shared style file
\texttt{indstyle.sty}, the bibliography \texttt{references.bib}, and the
\texttt{sections/} directory of section files, with every figure rendered as vector
TikZ inside the \texttt{mermaidfig} environment and every table set full width with
the ragged-right column types defined in the style file. There are no raster images,
no \texttt{\textbackslash includegraphics} calls, and no external image
dependencies, so the package compiles from its own sources without additional
assets. The bibliography is processed with the \texttt{ieeetr} style, and the
codified references throughout the package use the section symbol, for example 21
CFR \S 312.23, exactly as rendered by the style file. A reviewer can therefore
reproduce the compiled PDF directly from the repository at version v4.3.0 by
uploading the stage directory to Overleaf and compiling with pdfLaTeX, which closes
the loop from grounded source to verifiable, reproducible document
\cite{phase1guidance, daraxwhipple}.
